\documentclass[10pt,a4paper]{report}

% --- Essential Packages ---
\usepackage[utf8]{inputenc}
\usepackage[margin=0.75in,headheight=14pt]{geometry}
\usepackage{amsmath,amssymb}
\usepackage{xcolor}
\usepackage{tcolorbox}
\usepackage{listings}
\usepackage{fancyhdr}
\usepackage{titlesec}
\usepackage{booktabs}
\usepackage{enumitem}
\usepackage{hyperref}
\usepackage{bookmark}

% --- Color Palette ---
\definecolor{primary}{HTML}{0F172A}       % Deep Slate 900
\definecolor{secondary}{HTML}{1E40AF}     % Cobalt 800
\definecolor{accentteal}{HTML}{0D9488}    % Teal 600
\definecolor{codebg}{HTML}{F8FAFC}        % Off-white editor tint
\definecolor{termbg}{HTML}{0F172A}        % Dark terminal background
\definecolor{termfg}{HTML}{38BDF8}        % Cyan/Mint terminal text
\definecolor{theorybg}{HTML}{F0FDF4}      % Light emerald tint
\definecolor{theoryframe}{HTML}{16A34A}   % Emerald border

% --- Syntax Highlighting Colors ---
\definecolor{pykeyword}{HTML}{7C3AED}     % Violet bold
\definecolor{pystring}{HTML}{047857}      % Emerald
\definecolor{pycomment}{HTML}{64748B}     % Slate italic
\definecolor{pynumber}{HTML}{C2410C}      % Orange/Coral
\definecolor{pybuiltin}{HTML}{2563EB}     % Blue bold

% --- Listings Definitions ---
\lstdefinestyle{python_detailed}{
    backgroundcolor=\color{codebg},
    basicstyle=\ttfamily\scriptsize\color{primary},
    keywordstyle=\color{pykeyword}\bfseries,
    commentstyle=\color{pycomment}\itshape,
    stringstyle=\color{pystring},
    numberstyle=\tiny\color{pycomment},
    numbers=left,
    numbersep=6pt,
    breaklines=true,
    breakatwhitespace=false,
    showstringspaces=false,
    keepspaces=true,
    tabsize=4,
    morekeywords={with,as,True,False,None,self,append,round,range,len,abs,float,int}
}

\lstdefinestyle{terminal_output}{
    backgroundcolor=\color{termbg},
    basicstyle=\ttfamily\tiny\color{termfg},
    breaklines=true,
    breakatwhitespace=false,
    showstringspaces=false,
    keepspaces=true,
    numbers=none
}

% --- Chapter & Section Heading Styling ---
\titleformat{\chapter}[display]
  {\normalfont\huge\bfseries\color{primary}}
  {\filleft\Huge\textbf{\color{secondary}Chapter \thechapter}}
  {1ex}
  {\titlerule[2pt]\vspace{1ex}}
  [\vspace{1ex}\titlerule]

\titleformat{\section}
  {\normalfont\Large\bfseries\color{secondary}}
  {\thesection}{1em}{}

% --- Page Header & Footer ---
\pagestyle{fancy}
\fancyhf{}
\fancyhead[L]{\small\textbf{\color{secondary}Numerical Computing in Python}}
\fancyhead[R]{\small\nouppercase{\color{pycomment}\leftmark}}
\fancyfoot[L]{\scriptsize\textbf{\color{secondary}Code by Anushka}}
\fancyfoot[R]{\small\thepage}
\renewcommand{\headrulewidth}{0.6pt}
\renewcommand{\footrulewidth}{0.4pt}

% --- Hyperlink Setup ---
\hypersetup{
    colorlinks=true,
    linkcolor=secondary,
    urlcolor=accentteal,
    pdftitle={Numerical Computing Python Complete Handbook},
    pdfauthor={Code by Anushka},
    pdfsubject={Numerical Analysis and Computing Methods in Python}
}

\begin{document}

% ==========================================
% COVER PAGE
% ==========================================
\begin{titlepage}
\begin{center}
\vspace*{1.5cm}

{\Huge\bfseries\color{primary} NUMERICAL ANALYSIS \& COMPUTING}\\[0.4cm]
{\LARGE\bfseries\color{secondary} IN PYTHON}\\[0.8cm]

\textcolor{accentteal}{\rule{0.8\textwidth}{2.5pt}}\\[0.8cm]

{\Large\color{primary} A Complete Practical Handbook: Theoretical Formulations, Step-by-Step Procedural Implementations \& Iteration Convergence Tables}\\[1.5cm]

\begin{tcolorbox}[colback=codebg,colframe=secondary,width=0.88\textwidth,arc=4mm,boxrule=1.5pt]
\begin{center}
\textbf{\large\color{secondary} 20 Core Numerical Methods Included}\\[0.3cm]
\small\color{primary}
$\bullet$ Roots of Non-Linear Equations \quad $\bullet$ Systems of Linear Equations \\
$\bullet$ Interpolation \& Polynomials \quad $\bullet$ Numerical Quadrature \\
$\bullet$ Ordinary Differential Equations \quad $\bullet$ Eigenvalue Problems \& Fitting\\[0.2cm]
\textit{Featuring Clean Python 3 Syntax, Pure Procedural Functions, and Formatted Iteration Tables}
\end{center}
\end{tcolorbox}

\vfill

{\LARGE\textbf{\color{secondary}Code by Anushka}}\\[0.3cm]
{\small Semester 5 Mathematics \& Numerical Computing Laboratory Reference}\\[0.3cm]
{\small\today}

\end{center}
\end{titlepage}

% ==========================================
% TABLE OF CONTENTS
% ==========================================
\tableofcontents
\newpage
\chapter{Transcendental \& Polynomial Equations}

\section{Bisection Method}
\textbf{Category:} \textcolor{secondary}{Roots of Non-Linear Equations} \quad | \quad \textbf{Code by:} \textcolor{secondary}{Anushka} \quad | \quad \textbf{Source:} \texttt{\detokenize{01_bisection_detailed.py}}

The Bisection method is a robust bracketing method for finding real roots of a continuous function $f(x) = 0$. Given an initial bracket $[a, b]$ where $f(a) \cdot f(b) < 0$, Bolzano's theorem guarantees at least one root in $(a, b)$. In each iteration, the midpoint $c$ is calculated and the sign of $f(c)$ determines which sub-interval is retained.

\begin{tcolorbox}[colback=theorybg,colframe=theoryframe,title=\textbf{Mathematical Theory \& Formulation},boxrule=1pt,arc=2mm]
\[
c = \frac{a + b}{2.0}, \quad \text{Update: } b = c \text{ if } f(a) \cdot f(c) < 0, \quad a = c \text{ if } f(a) \cdot f(c) > 0
\]
\[
\text{Convergence: Linear}, \quad |e_{n+1}| \le \frac{1}{2}|e_n|, \quad \text{Error Bound: } \frac{b - a}{2^n}
\]
\end{tcolorbox}

\begin{tcolorbox}[colback=codebg,colframe=secondary,title=\textbf{Python Implementation (Step-by-Step Loops \& Pure Functions)},boxrule=1pt,arc=2mm]
\begin{lstlisting}[style=python_detailed]
# Code by Anushka

def bisection(f, a, b, tol=1e-5, max_iter=50):
    if f(a) * f(b) >= 0:
        print("Error: f(a) and f(b) must have opposite signs.")
        return None

    print(f"{'Iter':<6}{'a':<12}{'b':<12}{'c (Mid)':<14}{'f(c)':<12}")
    print("-" * 56)

    for i in range(1, max_iter + 1):
        c = (a + b) / 2.0
        fc = f(c)
        print(f"{i:<6}{a:<12.5f}{b:<12.5f}{c:<14.5f}{fc:<12.5f}")

        if abs(fc) < tol or (b - a) / 2.0 < tol:
            return c

        if f(a) * fc < 0:
            b = c
        else:
            a = c

    return (a + b) / 2.0


def f(x):
    return x**3 - x - 2


root = bisection(f, 1.0, 2.0)
print(f"\nRoot: {root:.5f}")

\end{lstlisting}
\end{tcolorbox}

\begin{tcolorbox}[colback=termbg,colframe=primary,title=\textbf{Program Execution Output \& Convergence Table},boxrule=1pt,arc=2mm]
\begin{lstlisting}[style=terminal_output]
Iter  a           b           c (Mid)       f(c)        
--------------------------------------------------------
1     1.00000     2.00000     1.50000       -0.12500    
2     1.50000     2.00000     1.75000       1.60938     
3     1.50000     1.75000     1.62500       0.66602     
4     1.50000     1.62500     1.56250       0.25220     
5     1.50000     1.56250     1.53125       0.05911     
6     1.50000     1.53125     1.51562       -0.03405    
7     1.51562     1.53125     1.52344       0.01225     
8     1.51562     1.52344     1.51953       -0.01097    
9     1.51953     1.52344     1.52148       0.00062     
10    1.51953     1.52148     1.52051       -0.00518    
11    1.52051     1.52148     1.52100       -0.00228    
12    1.52100     1.52148     1.52124       -0.00083    
13    1.52124     1.52148     1.52136       -0.00010    
14    1.52136     1.52148     1.52142       0.00026     
15    1.52136     1.52142     1.52139       0.00008     
16    1.52136     1.52139     1.52138       -0.00001    
17    1.52138     1.52139     1.52139       0.00003     

Root: 1.52139
\end{lstlisting}
\end{tcolorbox}

\newpage

\section{Regula-Falsi (False Position) Method}
\textbf{Category:} \textcolor{secondary}{Roots of Non-Linear Equations} \quad | \quad \textbf{Code by:} \textcolor{secondary}{Anushka} \quad | \quad \textbf{Source:} \texttt{\detokenize{02_regula_falsi_detailed.py}}

The Regula-Falsi (False Position) method combines the reliability of bracketing with the acceleration of linear interpolation. Instead of bisecting the interval geometrically, it connects $(a, f(a))$ and $(b, f(b))$ with a secant line and determines the root estimate $c$ at the line's $x$-intercept.

\begin{tcolorbox}[colback=theorybg,colframe=theoryframe,title=\textbf{Mathematical Theory \& Formulation},boxrule=1pt,arc=2mm]
\[
c = \frac{a f(b) - b f(a)}{f(b) - f(a)} = a - \frac{f(a)(b - a)}{f(b) - f(a)}
\]
\[
\text{Update: } b = c \text{ if } f(a) \cdot f(c) < 0, \quad a = c \text{ if } f(a) \cdot f(c) > 0
\]
\end{tcolorbox}

\begin{tcolorbox}[colback=codebg,colframe=secondary,title=\textbf{Python Implementation (Step-by-Step Loops \& Pure Functions)},boxrule=1pt,arc=2mm]
\begin{lstlisting}[style=python_detailed]
# Code by Anushka

def regula_falsi(f, a, b, tol=1e-5, max_iter=50):
    if f(a) * f(b) >= 0:
        print("Error: f(a) and f(b) must have opposite signs.")
        return None

    print(f"{'Iter':<6}{'a':<12}{'b':<12}{'c (False)':<14}{'f(c)':<12}")
    print("-" * 56)

    c = a
    for i in range(1, max_iter + 1):
        fa = f(a)
        fb = f(b)
        c = (a * fb - b * fa) / (fb - fa)
        fc = f(c)
        print(f"{i:<6}{a:<12.5f}{b:<12.5f}{c:<14.5f}{fc:<12.5f}")

        if abs(fc) < tol:
            return c

        if fa * fc < 0:
            b = c
        else:
            a = c

    return c


def f(x):
    return x**3 - 2 * x - 5


root = regula_falsi(f, 2.0, 3.0)
print(f"\nRoot: {root:.5f}")

\end{lstlisting}
\end{tcolorbox}

\begin{tcolorbox}[colback=termbg,colframe=primary,title=\textbf{Program Execution Output \& Convergence Table},boxrule=1pt,arc=2mm]
\begin{lstlisting}[style=terminal_output]
Iter  a           b           c (False)     f(c)        
--------------------------------------------------------
1     2.00000     3.00000     2.05882       -0.39080    
2     2.05882     3.00000     2.08126       -0.14720    
3     2.08126     3.00000     2.08964       -0.05468    
4     2.08964     3.00000     2.09274       -0.02020    
5     2.09274     3.00000     2.09388       -0.00745    
6     2.09388     3.00000     2.09431       -0.00275    
7     2.09431     3.00000     2.09446       -0.00101    
8     2.09446     3.00000     2.09452       -0.00037    
9     2.09452     3.00000     2.09454       -0.00014    
10    2.09454     3.00000     2.09455       -0.00005    
11    2.09455     3.00000     2.09455       -0.00002    
12    2.09455     3.00000     2.09455       -0.00001    

Root: 2.09455
\end{lstlisting}
\end{tcolorbox}

\newpage

\section{Newton-Raphson Method}
\textbf{Category:} \textcolor{secondary}{Roots of Non-Linear Equations} \quad | \quad \textbf{Code by:} \textcolor{secondary}{Anushka} \quad | \quad \textbf{Source:} \texttt{\detokenize{03_newton_raphson_detailed.py}}

The Newton-Raphson method is an open root-finding technique that utilizes the tangent line at the current estimate $x_k$. Starting from an initial guess $x_0$, the method converges with quadratic order near simple roots.

\begin{tcolorbox}[colback=theorybg,colframe=theoryframe,title=\textbf{Mathematical Theory \& Formulation},boxrule=1pt,arc=2mm]
\[
x_{k+1} = x_k - \frac{f(x_k)}{f'(x_k)}
\]
\[
\text{Order of Convergence: } p = 2 \quad (\text{Quadratic, requires } f'(x_k) \neq 0)
\]
\end{tcolorbox}

\begin{tcolorbox}[colback=codebg,colframe=secondary,title=\textbf{Python Implementation (Step-by-Step Loops \& Pure Functions)},boxrule=1pt,arc=2mm]
\begin{lstlisting}[style=python_detailed]
# Code by Anushka

def newton_raphson(f, df, x0, tol=1e-5, max_iter=50):
    print(f"{'Iter':<6}{'x0':<12}{'f(x0)':<14}{'df(x0)':<14}{'x1':<12}")
    print("-" * 58)

    for i in range(1, max_iter + 1):
        fx = f(x0)
        dfx = df(x0)
        if dfx == 0:
            print("Error: Derivative is zero.")
            return None

        x1 = x0 - (fx / dfx)
        print(f"{i:<6}{x0:<12.5f}{fx:<14.5f}{dfx:<14.5f}{x1:<12.5f}")

        if abs(x1 - x0) < tol or abs(f(x1)) < tol:
            return x1

        x0 = x1

    return x0


def f(x):
    return x**3 - 3 * x + 1


def df(x):
    return 3 * x**2 - 3


root = newton_raphson(f, df, 1.5)
print(f"\nRoot: {root:.5f}")

\end{lstlisting}
\end{tcolorbox}

\begin{tcolorbox}[colback=termbg,colframe=primary,title=\textbf{Program Execution Output \& Convergence Table},boxrule=1pt,arc=2mm]
\begin{lstlisting}[style=terminal_output]
Iter  x0          f(x0)         df(x0)        x1          
----------------------------------------------------------
1     1.50000     -0.12500      3.75000       1.53333     
2     1.53333     0.00504       4.05333       1.53209     

Root: 1.53209
\end{lstlisting}
\end{tcolorbox}

\newpage

\section{Secant Method}
\textbf{Category:} \textcolor{secondary}{Roots of Non-Linear Equations} \quad | \quad \textbf{Code by:} \textcolor{secondary}{Anushka} \quad | \quad \textbf{Source:} \texttt{\detokenize{04_secant_detailed.py}}

The Secant method approximates the analytical derivative in Newton-Raphson using a finite-difference slope between the two most recent estimates $x_{k-1}$ and $x_k$. This avoids calculating $f'(x)$ while retaining rapid superlinear convergence.

\begin{tcolorbox}[colback=theorybg,colframe=theoryframe,title=\textbf{Mathematical Theory \& Formulation},boxrule=1pt,arc=2mm]
\[
x_{k+1} = x_k - f(x_k) \cdot \frac{x_k - x_{k-1}}{f(x_k) - f(x_{k-1})}
\]
\[
\text{Order of Convergence: } p = \frac{1 + \sqrt{5}}{2} \approx 1.618 \quad (\text{Golden Ratio, Superlinear})
\]
\end{tcolorbox}

\begin{tcolorbox}[colback=codebg,colframe=secondary,title=\textbf{Python Implementation (Step-by-Step Loops \& Pure Functions)},boxrule=1pt,arc=2mm]
\begin{lstlisting}[style=python_detailed]
# Code by Anushka

def secant(f, x0, x1, tol=1e-5, max_iter=50):
    print(f"{'Iter':<6}{'x0':<12}{'x1':<12}{'x2':<14}{'f(x2)':<12}")
    print("-" * 56)

    for i in range(1, max_iter + 1):
        f0 = f(x0)
        f1 = f(x1)
        if f1 - f0 == 0:
            print("Error: Division by zero in Secant formula.")
            return None

        x2 = x1 - f1 * (x1 - x0) / (f1 - f0)
        f2 = f(x2)
        print(f"{i:<6}{x0:<12.5f}{x1:<12.5f}{x2:<14.5f}{f2:<12.5f}")

        if abs(x2 - x1) < tol or abs(f2) < tol:
            return x2

        x0 = x1
        x1 = x2

    return x1


def f(x):
    return x**2 - 4


root = secant(f, 1.0, 3.0)
print(f"\nRoot: {root:.5f}")

\end{lstlisting}
\end{tcolorbox}

\begin{tcolorbox}[colback=termbg,colframe=primary,title=\textbf{Program Execution Output \& Convergence Table},boxrule=1pt,arc=2mm]
\begin{lstlisting}[style=terminal_output]
Iter  x0          x1          x2            f(x2)       
--------------------------------------------------------
1     1.00000     3.00000     1.75000       -0.93750    
2     3.00000     1.75000     1.94737       -0.20776    
3     1.75000     1.94737     2.00356       0.01425     
4     1.94737     2.00356     1.99995       -0.00019    
5     2.00356     1.99995     2.00000       -0.00000    

Root: 2.00000
\end{lstlisting}
\end{tcolorbox}

\newpage

\section{Fixed-Point Iteration Method}
\textbf{Category:} \textcolor{secondary}{Roots of Non-Linear Equations} \quad | \quad \textbf{Code by:} \textcolor{secondary}{Anushka} \quad | \quad \textbf{Source:} \texttt{\detokenize{05_fixed_point_detailed.py}}

Fixed-point iteration rewrites the root problem $f(x) = 0$ as $x = g(x)$. Starting from $x_0$, successive estimates are generated by repeatedly evaluating $g(x)$. Convergence is guaranteed if $|g'(x)| < 1$ in an interval containing the fixed point.

\begin{tcolorbox}[colback=theorybg,colframe=theoryframe,title=\textbf{Mathematical Theory \& Formulation},boxrule=1pt,arc=2mm]
\[
x_{k+1} = g(x_k)
\]
\[
\text{Convergence Condition: } |g'(x^*)| < 1 \quad (\text{Contraction Mapping})
\]
\end{tcolorbox}

\begin{tcolorbox}[colback=codebg,colframe=secondary,title=\textbf{Python Implementation (Step-by-Step Loops \& Pure Functions)},boxrule=1pt,arc=2mm]
\begin{lstlisting}[style=python_detailed]
# Code by Anushka

import math

def fixed_point(g, x0, tol=1e-5, max_iter=50):
    print(f"{'Iter':<6}{'x0':<12}{'x1 = g(x0)':<16}{'Error':<12}")
    print("-" * 46)

    for i in range(1, max_iter + 1):
        x1 = g(x0)
        err = abs(x1 - x0)
        print(f"{i:<6}{x0:<12.5f}{x1:<16.5f}{err:<12.5f}")

        if err < tol:
            return x1

        x0 = x1

    return x0


def g(x):
    return math.cos(x)


root = fixed_point(g, 0.5)
print(f"\nFixed Point: {root:.5f}")

\end{lstlisting}
\end{tcolorbox}

\begin{tcolorbox}[colback=termbg,colframe=primary,title=\textbf{Program Execution Output \& Convergence Table},boxrule=1pt,arc=2mm]
\begin{lstlisting}[style=terminal_output]
Iter  x0          x1 = g(x0)      Error       
----------------------------------------------
1     0.50000     0.87758         0.37758     
2     0.87758     0.63901         0.23857     
3     0.63901     0.80269         0.16367     
4     0.80269     0.69478         0.10791     
5     0.69478     0.76820         0.07342     
6     0.76820     0.71917         0.04903     
7     0.71917     0.75236         0.03319     
8     0.75236     0.73008         0.02227     
9     0.73008     0.74512         0.01504     
10    0.74512     0.73501         0.01011     
11    0.73501     0.74183         0.00682     
12    0.74183     0.73724         0.00459     
13    0.73724     0.74033         0.00309     
14    0.74033     0.73825         0.00208     
15    0.73825     0.73965         0.00140     
16    0.73965     0.73870         0.00095     
17    0.73870     0.73934         0.00064     
18    0.73934     0.73891         0.00043     
19    0.73891     0.73920         0.00029     
20    0.73920     0.73901         0.00019     
21    0.73901     0.73914         0.00013     
22    0.73914     0.73905         0.00009     
23    0.73905     0.73911         0.00006     
24    0.73911     0.73907         0.00004     
25    0.73907     0.73910         0.00003     
26    0.73910     0.73908         0.00002     
27    0.73908     0.73909         0.00001     
28    0.73909     0.73908         0.00001     

Fixed Point: 0.73908
\end{lstlisting}
\end{tcolorbox}

\newpage

\chapter{Systems of Linear Algebraic Equations}

\section{Gauss Elimination Method}
\textbf{Category:} \textcolor{secondary}{Systems of Linear Algebraic Equations} \quad | \quad \textbf{Code by:} \textcolor{secondary}{Anushka} \quad | \quad \textbf{Source:} \texttt{\detokenize{06_gauss_elimination_detailed.py}}

Gauss Elimination is a direct method for solving a system of linear equations $A x = b$. It transforms the augmented matrix $[A \mid b]$ into upper-triangular form through forward row operations, followed by backward substitution to determine the unknowns.

\begin{tcolorbox}[colback=theorybg,colframe=theoryframe,title=\textbf{Mathematical Theory \& Formulation},boxrule=1pt,arc=2mm]
\[
\text{Forward Elimination: } m_{ik} = \frac{a_{ik}}{a_{kk}}, \quad R_i \leftarrow R_i - m_{ik} R_k \quad (i > k)
\]
\[
\text{Back Substitution: } x_i = \frac{1}{a_{ii}}\left(b_i - \sum_{j=i+1}^n a_{ij} x_j\right)
\]
\end{tcolorbox}

\begin{tcolorbox}[colback=codebg,colframe=secondary,title=\textbf{Python Implementation (Step-by-Step Loops \& Pure Functions)},boxrule=1pt,arc=2mm]
\begin{lstlisting}[style=python_detailed]
# Code by Anushka

def gauss_elimination(A, b):
    n = len(b)

    # Forward Elimination
    for i in range(n):
        for j in range(i + 1, n):
            ratio = A[j][i] / A[i][i]
            for k in range(i, n):
                A[j][k] = A[j][k] - ratio * A[i][k]
            b[j] = b[j] - ratio * b[i]

    # Back Substitution
    x = []
    for _ in range(n):
        x.append(0.0)

    for i in range(n - 1, -1, -1):
        sum_ax = 0.0
        for j in range(i + 1, n):
            sum_ax = sum_ax + A[i][j] * x[j]
        x[i] = (b[i] - sum_ax) / A[i][i]

    return x


A = [
    [2.0, 1.0, -1.0],
    [-3.0, -1.0, 2.0],
    [-2.0, 1.0, 2.0]
]
b = [8.0, -11.0, -3.0]

ans = gauss_elimination(A, b)
for i in range(len(ans)):
    print(f"x{i+1} = {ans[i]:.4f}")

\end{lstlisting}
\end{tcolorbox}

\begin{tcolorbox}[colback=termbg,colframe=primary,title=\textbf{Program Execution Output \& Convergence Table},boxrule=1pt,arc=2mm]
\begin{lstlisting}[style=terminal_output]
x1 = 2.0000
x2 = 3.0000
x3 = -1.0000
\end{lstlisting}
\end{tcolorbox}

\newpage

\section{Gauss-Jordan Elimination Method}
\textbf{Category:} \textcolor{secondary}{Systems of Linear Algebraic Equations} \quad | \quad \textbf{Code by:} \textcolor{secondary}{Anushka} \quad | \quad \textbf{Source:} \texttt{\detokenize{07_gauss_jordan_detailed.py}}

Gauss-Jordan elimination reduces the coefficient matrix directly into reduced row echelon form (identity matrix $I$) by performing row eliminations both below and above each pivot element. The solution appears directly in the augmented column.

\begin{tcolorbox}[colback=theorybg,colframe=theoryframe,title=\textbf{Mathematical Theory \& Formulation},boxrule=1pt,arc=2mm]
\[
\text{Normalize Pivot Row: } R_i \leftarrow \frac{R_i}{a_{ii}}
\]
\[
\text{Eliminate Other Rows: } R_k \leftarrow R_k - a_{ki} R_i \quad (\forall k \neq i)
\]
\[
\text{Result: } [A \mid b] \longrightarrow [I \mid x]
\]
\end{tcolorbox}

\begin{tcolorbox}[colback=codebg,colframe=secondary,title=\textbf{Python Implementation (Step-by-Step Loops \& Pure Functions)},boxrule=1pt,arc=2mm]
\begin{lstlisting}[style=python_detailed]
# Code by Anushka

def gauss_jordan(A, b):
    n = len(b)

    for i in range(n):
        pivot = A[i][i]

        # Normalize the pivot row
        for j in range(n):
            A[i][j] = A[i][j] / pivot
        b[i] = b[i] / pivot

        # Eliminate all other rows
        for k in range(n):
            if k != i:
                factor = A[k][i]
                for j in range(n):
                    A[k][j] = A[k][j] - factor * A[i][j]
                b[k] = b[k] - factor * b[i]

    return b


A = [
    [2.0, 3.0, 1.0],
    [1.0, 2.0, 3.0],
    [3.0, 1.0, 2.0]
]
b = [9.0, 6.0, 8.0]

ans = gauss_jordan(A, b)
for i in range(len(ans)):
    print(f"x{i+1} = {ans[i]:.4f}")

\end{lstlisting}
\end{tcolorbox}

\begin{tcolorbox}[colback=termbg,colframe=primary,title=\textbf{Program Execution Output \& Convergence Table},boxrule=1pt,arc=2mm]
\begin{lstlisting}[style=terminal_output]
x1 = 1.9444
x2 = 1.6111
x3 = 0.2778
\end{lstlisting}
\end{tcolorbox}

\newpage

\section{Gauss-Jacobi Iterative Method}
\textbf{Category:} \textcolor{secondary}{Systems of Linear Algebraic Equations} \quad | \quad \textbf{Code by:} \textcolor{secondary}{Anushka} \quad | \quad \textbf{Source:} \texttt{\detokenize{08_gauss_jacobi_detailed.py}}

The Jacobi method is an iterative technique for solving $A x = b$. In each step, each unknown $x_i$ is updated solely using the values from the previous iteration step $k$. It converges for strictly diagonally dominant matrices.

\begin{tcolorbox}[colback=theorybg,colframe=theoryframe,title=\textbf{Mathematical Theory \& Formulation},boxrule=1pt,arc=2mm]
\[
x_i^{(k+1)} = \frac{1}{a_{ii}}\left(b_i - \sum_{j \neq i} a_{ij} x_j^{(k)}\right) \quad (i = 1, 2, \dots, n)
\]
\[
\text{Strict Diagonal Dominance: } |a_{ii}| > \sum_{j \neq i} |a_{ij}|
\]
\end{tcolorbox}

\begin{tcolorbox}[colback=codebg,colframe=secondary,title=\textbf{Python Implementation (Step-by-Step Loops \& Pure Functions)},boxrule=1pt,arc=2mm]
\begin{lstlisting}[style=python_detailed]
# Code by Anushka

def gauss_jacobi(A, b, tol=1e-5, max_iter=50):
    n = len(b)
    x = []
    for _ in range(n):
        x.append(0.0)

    # Print Table Header
    header = f"{'Iter':<6}"
    for i in range(n):
        header += f"{'x' + str(i+1):<12}"
    print(header)
    print("-" * (6 + 12 * n))

    for step in range(1, max_iter + 1):
        x_new = []
        for i in range(n):
            sum_other = 0.0
            for j in range(n):
                if j != i:
                    sum_other += A[i][j] * x[j]
            xi_val = (b[i] - sum_other) / A[i][i]
            x_new.append(xi_val)

        # Print current iteration values
        row_str = f"{step:<6}"
        for val in x_new:
            row_str += f"{val:<12.5f}"
        print(row_str)

        # Check convergence
        max_diff = 0.0
        for i in range(n):
            diff = abs(x_new[i] - x[i])
            if diff > max_diff:
                max_diff = diff

        if max_diff < tol:
            return x_new

        x = x_new

    return x


A = [
    [10.0, -1.0, 2.0],
    [-1.0, 11.0, -1.0],
    [2.0, -1.0, 10.0]
]
b = [6.0, 25.0, -11.0]

sol = gauss_jacobi(A, b)
rounded_sol = []
for v in sol:
    rounded_sol.append(round(v, 4))
print("\nSolution:", rounded_sol)

\end{lstlisting}
\end{tcolorbox}

\begin{tcolorbox}[colback=termbg,colframe=primary,title=\textbf{Program Execution Output \& Convergence Table},boxrule=1pt,arc=2mm]
\begin{lstlisting}[style=terminal_output]
Iter  x1          x2          x3          
------------------------------------------
1     0.60000     2.27273     -1.10000    
2     1.04727     2.22727     -0.99273    
3     1.02127     2.27769     -1.08673    
4     1.04511     2.26678     -1.07649    
5     1.04197     2.26988     -1.08235    
6     1.04346     2.26906     -1.08141    
7     1.04319     2.26928     -1.08179    
8     1.04328     2.26922     -1.08171    
9     1.04326     2.26923     -1.08174    
10    1.04327     2.26923     -1.08173    

Solution: [1.0433, 2.2692, -1.0817]
\end{lstlisting}
\end{tcolorbox}

\newpage

\section{Gauss-Seidel Iterative Method}
\textbf{Category:} \textcolor{secondary}{Systems of Linear Algebraic Equations} \quad | \quad \textbf{Code by:} \textcolor{secondary}{Anushka} \quad | \quad \textbf{Source:} \texttt{\detokenize{09_gauss_seidel_detailed.py}}

The Gauss-Seidel method accelerates Jacobi iteration by immediately utilizing the newly updated components $x_1^{(k+1)}, \dots, x_{i-1}^{(k+1)}$ as soon as they become available during the current iteration step.

\begin{tcolorbox}[colback=theorybg,colframe=theoryframe,title=\textbf{Mathematical Theory \& Formulation},boxrule=1pt,arc=2mm]
\[
x_i^{(k+1)} = \frac{1}{a_{ii}}\left(b_i - \sum_{j < i} a_{ij} x_j^{(k+1)} - \sum_{j > i} a_{ij} x_j^{(k)}\right)
\]
\[
\text{Convergence Rate: Typically converges in roughly half the iterations of Jacobi.}
\]
\end{tcolorbox}

\begin{tcolorbox}[colback=codebg,colframe=secondary,title=\textbf{Python Implementation (Step-by-Step Loops \& Pure Functions)},boxrule=1pt,arc=2mm]
\begin{lstlisting}[style=python_detailed]
# Code by Anushka

def gauss_seidel(A, b, tol=1e-5, max_iter=50):
    n = len(b)
    x = []
    for _ in range(n):
        x.append(0.0)

    header = f"{'Iter':<6}"
    for i in range(n):
        header += f"{'x' + str(i+1):<12}"
    print(header)
    print("-" * (6 + 12 * n))

    for step in range(1, max_iter + 1):
        max_diff = 0.0

        for i in range(n):
            sum_other = 0.0
            for j in range(n):
                if j != i:
                    sum_other += A[i][j] * x[j]

            x_new_val = (b[i] - sum_other) / A[i][i]
            diff = abs(x_new_val - x[i])
            if diff > max_diff:
                max_diff = diff
            x[i] = x_new_val

        row_str = f"{step:<6}"
        for val in x:
            row_str += f"{val:<12.5f}"
        print(row_str)

        if max_diff < tol:
            return x

    return x


A = [
    [10.0, -1.0, 2.0],
    [-1.0, 11.0, -1.0],
    [2.0, -1.0, 10.0]
]
b = [6.0, 25.0, -11.0]

sol = gauss_seidel(A, b)
rounded_sol = []
for v in sol:
    rounded_sol.append(round(v, 4))
print("\nSolution:", rounded_sol)

\end{lstlisting}
\end{tcolorbox}

\begin{tcolorbox}[colback=termbg,colframe=primary,title=\textbf{Program Execution Output \& Convergence Table},boxrule=1pt,arc=2mm]
\begin{lstlisting}[style=terminal_output]
Iter  x1          x2          x3          
------------------------------------------
1     0.60000     2.32727     -0.98727    
2     1.03018     2.27663     -1.07837    
3     1.04334     2.26954     -1.08171    
4     1.04330     2.26923     -1.08174    
5     1.04327     2.26923     -1.08173    
6     1.04327     2.26923     -1.08173    

Solution: [1.0433, 2.2692, -1.0817]
\end{lstlisting}
\end{tcolorbox}

\newpage

\chapter{Interpolation \& Polynomial Approximation}

\section{Lagrange Interpolation}
\textbf{Category:} \textcolor{secondary}{Interpolation \& Polynomial Approximation} \quad | \quad \textbf{Code by:} \textcolor{secondary}{Anushka} \quad | \quad \textbf{Source:} \texttt{\detokenize{10_lagrange_detailed.py}}

Lagrange interpolation constructs the unique polynomial of degree at most $n$ that passes through $n+1$ given points $(x_i, y_i)$. The formulation does not require the data nodes to be evenly spaced.

\begin{tcolorbox}[colback=theorybg,colframe=theoryframe,title=\textbf{Mathematical Theory \& Formulation},boxrule=1pt,arc=2mm]
\[
P_n(x) = \sum_{i=0}^n y_i \ell_i(x), \quad \text{where } \ell_i(x) = \prod_{\substack{j=0 \\ j \neq i}}^n \frac{x - x_j}{x_i - x_j}
\]
\end{tcolorbox}

\begin{tcolorbox}[colback=codebg,colframe=secondary,title=\textbf{Python Implementation (Step-by-Step Loops \& Pure Functions)},boxrule=1pt,arc=2mm]
\begin{lstlisting}[style=python_detailed]
# Code by Anushka

def lagrange(x_vals, y_vals, target_x):
    n = len(x_vals)
    print(f"{'i':<6}{'Basis L_i(x)':<16}{'y_i':<10}{'Term':<12}")
    print("-" * 44)

    total = 0.0
    for i in range(n):
        # Calculate Lagrange basis polynomial L_i(x) using a loop
        L_i = 1.0
        for j in range(n):
            if i != j:
                L_i = L_i * (target_x - x_vals[j]) / (x_vals[i] - x_vals[j])

        term = y_vals[i] * L_i
        total += term
        print(f"{i:<6}{L_i:<16.5f}{y_vals[i]:<10.2f}{term:<12.5f}")

    return total


x_pts = [5.0, 6.0, 9.0, 11.0]
y_pts = [12.0, 13.0, 14.0, 16.0]

ans = lagrange(x_pts, y_pts, 10.0)
print(f"\nInterpolated value: {ans:.5f}")

\end{lstlisting}
\end{tcolorbox}

\begin{tcolorbox}[colback=termbg,colframe=primary,title=\textbf{Program Execution Output \& Convergence Table},boxrule=1pt,arc=2mm]
\begin{lstlisting}[style=terminal_output]
i     Basis L_i(x)    y_i       Term        
--------------------------------------------
0     0.16667         12.00     2.00000     
1     -0.33333        13.00     -4.33333    
2     0.83333         14.00     11.66667    
3     0.33333         16.00     5.33333     

Interpolated value: 14.66667
\end{lstlisting}
\end{tcolorbox}

\newpage

\section{Newton's Divided Difference Interpolation}
\textbf{Category:} \textcolor{secondary}{Interpolation \& Polynomial Approximation} \quad | \quad \textbf{Code by:} \textcolor{secondary}{Anushka} \quad | \quad \textbf{Source:} \texttt{\detokenize{11_newton_divided_diff_detailed.py}}

Newton's Divided Difference method constructs an interpolating polynomial using divided differences. It is versatile for unequally spaced data points and allows easy addition of new data points without recomputing the entire table.

\begin{tcolorbox}[colback=theorybg,colframe=theoryframe,title=\textbf{Mathematical Theory \& Formulation},boxrule=1pt,arc=2mm]
\[
f[x_i, \dots, x_{i+k}] = \frac{f[x_{i+1}, \dots, x_{i+k}] - f[x_i, \dots, x_{i+k-1}]}{x_{i+k} - x_i}
\]
\[
P_n(x) = f[x_0] + \sum_{k=1}^n f[x_0, \dots, x_k] \prod_{j=0}^{k-1} (x - x_j)
\]
\end{tcolorbox}

\begin{tcolorbox}[colback=codebg,colframe=secondary,title=\textbf{Python Implementation (Step-by-Step Loops \& Pure Functions)},boxrule=1pt,arc=2mm]
\begin{lstlisting}[style=python_detailed]
# Code by Anushka

def newton_divided_diff(x, y, target_x):
    n = len(x)

    # Initialize 2D table with zeros using loops
    table = []
    for i in range(n):
        row = []
        for j in range(n):
            row.append(0.0)
        table.append(row)

    # First column is y values
    for i in range(n):
        table[i][0] = y[i]

    # Compute divided differences
    for j in range(1, n):
        for i in range(n - j):
            table[i][j] = (table[i + 1][j - 1] - table[i][j - 1]) / (x[i + j] - x[i])

    print("Divided Difference Table:")
    for i in range(n):
        row_str = f"{x[i]:<6.1f}"
        for j in range(n - i):
            row_str += f"{table[i][j]:<12.4f}"
        print(row_str)

    # Evaluate the polynomial at target_x
    res = table[0][0]
    prod_term = 1.0
    for j in range(1, n):
        prod_term = prod_term * (target_x - x[j - 1])
        res += table[0][j] * prod_term

    return res


x_pts = [4.0, 5.0, 7.0, 10.0, 11.0]
y_pts = [48.0, 100.0, 294.0, 900.0, 1210.0]

ans = newton_divided_diff(x_pts, y_pts, 6.0)
print(f"\nInterpolated Value: {ans:.5f}")

\end{lstlisting}
\end{tcolorbox}

\begin{tcolorbox}[colback=termbg,colframe=primary,title=\textbf{Program Execution Output \& Convergence Table},boxrule=1pt,arc=2mm]
\begin{lstlisting}[style=terminal_output]
Divided Difference Table:
4.0   48.0000     52.0000     15.0000     1.0000      0.0000      
5.0   100.0000    97.0000     21.0000     1.0000      
7.0   294.0000    202.0000    27.0000     
10.0  900.0000    310.0000    
11.0  1210.0000   

Interpolated Value: 180.00000
\end{lstlisting}
\end{tcolorbox}

\newpage

\section{Newton's Forward Difference Interpolation}
\textbf{Category:} \textcolor{secondary}{Interpolation \& Polynomial Approximation} \quad | \quad \textbf{Code by:} \textcolor{secondary}{Anushka} \quad | \quad \textbf{Source:} \texttt{\detokenize{12_newton_forward_detailed.py}}

Newton's Forward Difference formula is designed for tabulated data where nodes are equally spaced ($x_i = x_0 + i h$). It is ideal for interpolating values near the start of the table.

\begin{tcolorbox}[colback=theorybg,colframe=theoryframe,title=\textbf{Mathematical Theory \& Formulation},boxrule=1pt,arc=2mm]
\[
u = \frac{x - x_0}{h}, \quad \Delta y_i = y_{i+1} - y_i
\]
\[
P_n(x) = y_0 + u \Delta y_0 + \frac{u(u-1)}{2!} \Delta^2 y_0 + \dots + \frac{u(u-1)\dots(u-n+1)}{n!} \Delta^n y_0
\]
\end{tcolorbox}

\begin{tcolorbox}[colback=codebg,colframe=secondary,title=\textbf{Python Implementation (Step-by-Step Loops \& Pure Functions)},boxrule=1pt,arc=2mm]
\begin{lstlisting}[style=python_detailed]
# Code by Anushka

import math

def newton_forward(x, y, target_x):
    n = len(x)
    h = x[1] - x[0]
    u = (target_x - x[0]) / h

    # Initialize forward difference table using loops
    diff = []
    for i in range(n):
        row = []
        for j in range(n):
            row.append(0.0)
        diff.append(row)

    # Row 0 contains y values
    for i in range(n):
        diff[0][i] = y[i]

    # Compute difference table
    for order in range(1, n):
        for i in range(n - order):
            diff[order][i] = diff[order - 1][i + 1] - diff[order - 1][i]

    print("Forward Difference Table:")
    for i in range(n):
        row_str = f"{x[i]:<6.1f}"
        for j in range(n - i):
            row_str += f"{diff[j][i]:<12.4f}"
        print(row_str)

    res = diff[0][0]
    u_term = 1.0
    for order in range(1, n):
        u_term = u_term * (u - (order - 1))
        res += (u_term * diff[order][0]) / math.factorial(order)

    return res


x_pts = [10.0, 20.0, 30.0, 40.0, 50.0]
y_pts = [0.1736, 0.3420, 0.5000, 0.6428, 0.7660]

ans = newton_forward(x_pts, y_pts, 15.0)
print(f"\nInterpolated Value: {ans:.5f}")

\end{lstlisting}
\end{tcolorbox}

\begin{tcolorbox}[colback=termbg,colframe=primary,title=\textbf{Program Execution Output \& Convergence Table},boxrule=1pt,arc=2mm]
\begin{lstlisting}[style=terminal_output]
Forward Difference Table:
10.0  0.1736      0.1684      -0.0104     -0.0048     0.0004      
20.0  0.3420      0.1580      -0.0152     -0.0044     
30.0  0.5000      0.1428      -0.0196     
40.0  0.6428      0.1232      
50.0  0.7660      

Interpolated Value: 0.25878
\end{lstlisting}
\end{tcolorbox}

\newpage

\chapter{Numerical Quadrature \& Integration}

\section{Trapezoidal Rule of Numerical Integration}
\textbf{Category:} \textcolor{secondary}{Numerical Quadrature \& Integration} \quad | \quad \textbf{Code by:} \textcolor{secondary}{Anushka} \quad | \quad \textbf{Source:} \texttt{\detokenize{13_trapezoidal_detailed.py}}

The composite Trapezoidal rule evaluates $\int_a^b f(x)\,dx$ by approximating the integrand with linear segments across $n$ uniform subintervals of width $h = \frac{b-a}{n}$.

\begin{tcolorbox}[colback=theorybg,colframe=theoryframe,title=\textbf{Mathematical Theory \& Formulation},boxrule=1pt,arc=2mm]
\[
I \approx \frac{h}{2}\left[f(a) + 2 \sum_{i=1}^{n-1} f(x_i) + f(b)\right], \quad h = \frac{b - a}{n}
\]
\[
\text{Global Truncation Error: } E_T = -\frac{(b - a) h^2}{12} f''(\xi) \quad (\mathcal{O}(h^2) \text{ Accuracy})
\]
\end{tcolorbox}

\begin{tcolorbox}[colback=codebg,colframe=secondary,title=\textbf{Python Implementation (Step-by-Step Loops \& Pure Functions)},boxrule=1pt,arc=2mm]
\begin{lstlisting}[style=python_detailed]
# Code by Anushka

def trapezoidal(f, a, b, n=6):
    h = (b - a) / float(n)

    x_pts = []
    for i in range(n + 1):
        x_pts.append(a + i * h)

    y_pts = []
    for x in x_pts:
        y_pts.append(f(x))

    print(f"{'i':<6}{'x_i':<12}{'f(x_i)':<14}")
    print("-" * 32)
    for i in range(len(x_pts)):
        print(f"{i:<6}{x_pts[i]:<12.4f}{y_pts[i]:<14.5f}")

    sum_middle = 0.0
    for i in range(1, n):
        sum_middle += y_pts[i]

    integral = (h / 2.0) * (y_pts[0] + 2.0 * sum_middle + y_pts[n])
    return integral


def f(x):
    return 1.0 / (1.0 + x**2)


ans = trapezoidal(f, 0.0, 1.0, n=6)
print(f"\nIntegral: {ans:.5f}")

\end{lstlisting}
\end{tcolorbox}

\begin{tcolorbox}[colback=termbg,colframe=primary,title=\textbf{Program Execution Output \& Convergence Table},boxrule=1pt,arc=2mm]
\begin{lstlisting}[style=terminal_output]
i     x_i         f(x_i)        
--------------------------------
0     0.0000      1.00000       
1     0.1667      0.97297       
2     0.3333      0.90000       
3     0.5000      0.80000       
4     0.6667      0.69231       
5     0.8333      0.59016       
6     1.0000      0.50000       

Integral: 0.78424
\end{lstlisting}
\end{tcolorbox}

\newpage

\section{Simpson's 1/3 Rule of Numerical Integration}
\textbf{Category:} \textcolor{secondary}{Numerical Quadrature \& Integration} \quad | \quad \textbf{Code by:} \textcolor{secondary}{Anushka} \quad | \quad \textbf{Source:} \texttt{\detokenize{14_simpson_1_3_detailed.py}}

Simpson's 1/3 rule fits quadratic parabolas over adjacent pairs of subintervals. The total number of intervals $n$ must be an even integer. It achieves fourth-order accuracy.

\begin{tcolorbox}[colback=theorybg,colframe=theoryframe,title=\textbf{Mathematical Theory \& Formulation},boxrule=1pt,arc=2mm]
\[
I \approx \frac{h}{3}\left[f(a) + f(b) + 4 \sum_{i \text{ odd}} f(x_i) + 2 \sum_{i \text{ even}} f(x_i)\right], \quad (n \text{ must be even})
\]
\[
\text{Global Truncation Error: } E_S = -\frac{(b - a) h^4}{180} f^{(4)}(\xi) \quad (\mathcal{O}(h^4) \text{ Accuracy})
\]
\end{tcolorbox}

\begin{tcolorbox}[colback=codebg,colframe=secondary,title=\textbf{Python Implementation (Step-by-Step Loops \& Pure Functions)},boxrule=1pt,arc=2mm]
\begin{lstlisting}[style=python_detailed]
# Code by Anushka

def simpson_1_3(f, a, b, n=6):
    if n % 2 != 0:
        n += 1

    h = (b - a) / float(n)

    x_pts = []
    for i in range(n + 1):
        x_pts.append(a + i * h)

    y_pts = []
    for x in x_pts:
        y_pts.append(f(x))

    print(f"{'i':<6}{'x_i':<12}{'f(x_i)':<14}{'Mult':<6}")
    print("-" * 38)
    for i in range(len(x_pts)):
        if i == 0 or i == n:
            mult = 1
        elif i % 2 != 0:
            mult = 4
        else:
            mult = 2
        print(f"{i:<6}{x_pts[i]:<12.4f}{y_pts[i]:<14.5f}{mult:<6}")

    sum_odd = 0.0
    for i in range(1, n, 2):
        sum_odd += y_pts[i]

    sum_even = 0.0
    for i in range(2, n, 2):
        sum_even += y_pts[i]

    integral = (h / 3.0) * (y_pts[0] + y_pts[n] + 4.0 * sum_odd + 2.0 * sum_even)
    return integral


def f(x):
    return 1.0 / (1.0 + x**2)


ans = simpson_1_3(f, 0.0, 1.0, n=6)
print(f"\nIntegral: {ans:.5f}")

\end{lstlisting}
\end{tcolorbox}

\begin{tcolorbox}[colback=termbg,colframe=primary,title=\textbf{Program Execution Output \& Convergence Table},boxrule=1pt,arc=2mm]
\begin{lstlisting}[style=terminal_output]
i     x_i         f(x_i)        Mult  
--------------------------------------
0     0.0000      1.00000       1     
1     0.1667      0.97297       4     
2     0.3333      0.90000       2     
3     0.5000      0.80000       4     
4     0.6667      0.69231       2     
5     0.8333      0.59016       4     
6     1.0000      0.50000       1     

Integral: 0.78540
\end{lstlisting}
\end{tcolorbox}

\newpage

\section{Simpson's 3/8 Rule of Numerical Integration}
\textbf{Category:} \textcolor{secondary}{Numerical Quadrature \& Integration} \quad | \quad \textbf{Code by:} \textcolor{secondary}{Anushka} \quad | \quad \textbf{Source:} \texttt{\detokenize{15_simpson_3_8_detailed.py}}

Simpson's 3/8 rule fits cubic polynomials across groups of three consecutive subintervals. The total number of subintervals $n$ must be a multiple of 3.

\begin{tcolorbox}[colback=theorybg,colframe=theoryframe,title=\textbf{Mathematical Theory \& Formulation},boxrule=1pt,arc=2mm]
\[
I \approx \frac{3h}{8}\left[f(a) + f(b) + 3 \sum_{i \nmid 3} f(x_i) + 2 \sum_{i \mid 3} f(x_i)\right], \quad (n \in 3\mathbb{Z})
\]
\[
\text{Global Truncation Error: } E_{3/8} = -\frac{(b - a) h^4}{80} f^{(4)}(\xi) \quad (\mathcal{O}(h^4) \text{ Accuracy})
\]
\end{tcolorbox}

\begin{tcolorbox}[colback=codebg,colframe=secondary,title=\textbf{Python Implementation (Step-by-Step Loops \& Pure Functions)},boxrule=1pt,arc=2mm]
\begin{lstlisting}[style=python_detailed]
# Code by Anushka

def simpson_3_8(f, a, b, n=6):
    while n % 3 != 0:
        n += 1

    h = (b - a) / float(n)

    x_pts = []
    for i in range(n + 1):
        x_pts.append(a + i * h)

    y_pts = []
    for x in x_pts:
        y_pts.append(f(x))

    print(f"{'i':<6}{'x_i':<12}{'f(x_i)':<14}{'Mult':<6}")
    print("-" * 38)
    for i in range(len(x_pts)):
        if i == 0 or i == n:
            mult = 1
        elif i % 3 == 0:
            mult = 2
        else:
            mult = 3
        print(f"{i:<6}{x_pts[i]:<12.4f}{y_pts[i]:<14.5f}{mult:<6}")

    sum_middle = 0.0
    for i in range(1, n):
        if i % 3 == 0:
            sum_middle += 2.0 * y_pts[i]
        else:
            sum_middle += 3.0 * y_pts[i]

    integral = (3.0 * h / 8.0) * (y_pts[0] + y_pts[n] + sum_middle)
    return integral


def f(x):
    return 1.0 / (1.0 + x**2)


ans = simpson_3_8(f, 0.0, 1.0, n=6)
print(f"\nIntegral: {ans:.5f}")

\end{lstlisting}
\end{tcolorbox}

\begin{tcolorbox}[colback=termbg,colframe=primary,title=\textbf{Program Execution Output \& Convergence Table},boxrule=1pt,arc=2mm]
\begin{lstlisting}[style=terminal_output]
i     x_i         f(x_i)        Mult  
--------------------------------------
0     0.0000      1.00000       1     
1     0.1667      0.97297       3     
2     0.3333      0.90000       3     
3     0.5000      0.80000       2     
4     0.6667      0.69231       3     
5     0.8333      0.59016       3     
6     1.0000      0.50000       1     

Integral: 0.78540
\end{lstlisting}
\end{tcolorbox}

\newpage

\chapter{Ordinary Differential Equations (ODEs)}

\section{Euler's Method for ODEs}
\textbf{Category:} \textcolor{secondary}{Ordinary Differential Equations (ODEs)} \quad | \quad \textbf{Code by:} \textcolor{secondary}{Anushka} \quad | \quad \textbf{Source:} \texttt{\detokenize{16_euler_detailed.py}}

Euler's method is the classical first-order explicit single-step method for solving initial value problems $\frac{dy}{dx} = f(x, y)$ with $y(x_0) = y_0$.

\begin{tcolorbox}[colback=theorybg,colframe=theoryframe,title=\textbf{Mathematical Theory \& Formulation},boxrule=1pt,arc=2mm]
\[
y_{n+1} = y_n + h \cdot f(x_n, y_n)
\]
\[
\text{Global Error Order: } \mathcal{O}(h) \quad (\text{First-Order Accurate})
\]
\end{tcolorbox}

\begin{tcolorbox}[colback=codebg,colframe=secondary,title=\textbf{Python Implementation (Step-by-Step Loops \& Pure Functions)},boxrule=1pt,arc=2mm]
\begin{lstlisting}[style=python_detailed]
# Code by Anushka

def euler(f, x0, y0, x_end, h=0.05):
    x = x0
    y = y0
    step = 0

    print(f"{'Step':<6}{'x':<10}{'y':<12}{'Slope':<12}")
    print("-" * 40)

    while x < x_end - 1e-9:
        slope = f(x, y)
        print(f"{step:<6}{x:<10.3f}{y:<12.5f}{slope:<12.5f}")
        y = y + h * slope
        x = x + h
        step = step + 1

    print(f"{step:<6}{x:<10.3f}{y:<12.5f}")
    return y


def f(x, y):
    return x + y


ans = euler(f, 0.0, 1.0, 0.2, h=0.05)
print(f"\ny(0.2): {ans:.5f}")

\end{lstlisting}
\end{tcolorbox}

\begin{tcolorbox}[colback=termbg,colframe=primary,title=\textbf{Program Execution Output \& Convergence Table},boxrule=1pt,arc=2mm]
\begin{lstlisting}[style=terminal_output]
Step  x         y           Slope       
----------------------------------------
0     0.000     1.00000     1.00000     
1     0.050     1.05000     1.10000     
2     0.100     1.10500     1.20500     
3     0.150     1.16525     1.31525     
4     0.200     1.23101     

y(0.2): 1.23101
\end{lstlisting}
\end{tcolorbox}

\newpage

\section{Modified Euler's (Heun's) Method}
\textbf{Category:} \textcolor{secondary}{Ordinary Differential Equations (ODEs)} \quad | \quad \textbf{Code by:} \textcolor{secondary}{Anushka} \quad | \quad \textbf{Source:} \texttt{\detokenize{17_modified_euler_detailed.py}}

Heun's (Modified Euler) method improves upon Euler's method by implementing a predictor-corrector scheme that averages slopes at both the start and predicted end of the step.

\begin{tcolorbox}[colback=theorybg,colframe=theoryframe,title=\textbf{Mathematical Theory \& Formulation},boxrule=1pt,arc=2mm]
\[
\text{Predictor: } \tilde{y}_{n+1} = y_n + h \cdot f(x_n, y_n)
\]
\[
\text{Corrector: } y_{n+1} = y_n + \frac{h}{2}\left[f(x_n, y_n) + f(x_{n+1}, \tilde{y}_{n+1})\right] \quad (\mathcal{O}(h^2) \text{ Accuracy})
\]
\end{tcolorbox}

\begin{tcolorbox}[colback=codebg,colframe=secondary,title=\textbf{Python Implementation (Step-by-Step Loops \& Pure Functions)},boxrule=1pt,arc=2mm]
\begin{lstlisting}[style=python_detailed]
# Code by Anushka

def modified_euler(f, x0, y0, x_end, h=0.05):
    x = x0
    y = y0
    step = 0

    print(f"{'Step':<6}{'x':<10}{'y':<12}{'y_predict':<14}")
    print("-" * 42)

    while x < x_end - 1e-9:
        k1 = f(x, y)
        y_predict = y + h * k1
        print(f"{step:<6}{x:<10.3f}{y:<12.5f}{y_predict:<14.5f}")
        k2 = f(x + h, y_predict)
        y = y + (h / 2.0) * (k1 + k2)
        x = x + h
        step = step + 1

    print(f"{step:<6}{x:<10.3f}{y:<12.5f}")
    return y


def f(x, y):
    return x + y


ans = modified_euler(f, 0.0, 1.0, 0.2, h=0.05)
print(f"\ny(0.2): {ans:.5f}")

\end{lstlisting}
\end{tcolorbox}

\begin{tcolorbox}[colback=termbg,colframe=primary,title=\textbf{Program Execution Output \& Convergence Table},boxrule=1pt,arc=2mm]
\begin{lstlisting}[style=terminal_output]
Step  x         y           y_predict     
------------------------------------------
0     0.000     1.00000     1.05000       
1     0.050     1.05250     1.10763       
2     0.100     1.11025     1.17077       
3     0.150     1.17353     1.23971       
4     0.200     1.24261     

y(0.2): 1.24261
\end{lstlisting}
\end{tcolorbox}

\newpage

\section{Runge-Kutta 4th Order Method (RK4)}
\textbf{Category:} \textcolor{secondary}{Ordinary Differential Equations (ODEs)} \quad | \quad \textbf{Code by:} \textcolor{secondary}{Anushka} \quad | \quad \textbf{Source:} \texttt{\detokenize{18_rk4_detailed.py}}

The classical Runge-Kutta 4th Order (RK4) method evaluates four intermediate slopes per step and combines them with Simpson-type weighting to achieve fourth-order accuracy without evaluating higher derivatives.

\begin{tcolorbox}[colback=theorybg,colframe=theoryframe,title=\textbf{Mathematical Theory \& Formulation},boxrule=1pt,arc=2mm]
\[
k_1 = h f(x_n, y_n), \quad k_2 = h f\left(x_n + \frac{h}{2}, y_n + \frac{k_1}{2}\right)
\]
\[
k_3 = h f\left(x_n + \frac{h}{2}, y_n + \frac{k_2}{2}\right), \quad k_4 = h f(x_n + h, y_n + k_3)
\]
\[
y_{n+1} = y_n + \frac{1}{6}(k_1 + 2k_2 + 2k_3 + k_4) \quad (\mathcal{O}(h^4) \text{ Accuracy})
\]
\end{tcolorbox}

\begin{tcolorbox}[colback=codebg,colframe=secondary,title=\textbf{Python Implementation (Step-by-Step Loops \& Pure Functions)},boxrule=1pt,arc=2mm]
\begin{lstlisting}[style=python_detailed]
# Code by Anushka

def rk4(f, x0, y0, x_end, h=0.05):
    x = x0
    y = y0
    step = 0

    print(f"{'Step':<6}{'x':<8}{'y':<10}{'k1':<9}{'k2':<9}{'k3':<9}{'k4':<9}")
    print("-" * 60)

    while x < x_end - 1e-9:
        k1 = h * f(x, y)
        k2 = h * f(x + h / 2.0, y + k1 / 2.0)
        k3 = h * f(x + h / 2.0, y + k2 / 2.0)
        k4 = h * f(x + h, y + k3)
        print(f"{step:<6}{x:<8.2f}{y:<10.4f}{k1:<9.4f}{k2:<9.4f}{k3:<9.4f}{k4:<9.4f}")
        y = y + (k1 + 2.0 * k2 + 2.0 * k3 + k4) / 6.0
        x = x + h
        step = step + 1

    print(f"{step:<6}{x:<8.2f}{y:<10.4f}")
    return y


def f(x, y):
    return x + y


ans = rk4(f, 0.0, 1.0, 0.2, h=0.05)
print(f"\ny(0.2): {ans:.5f}")

\end{lstlisting}
\end{tcolorbox}

\begin{tcolorbox}[colback=termbg,colframe=primary,title=\textbf{Program Execution Output \& Convergence Table},boxrule=1pt,arc=2mm]
\begin{lstlisting}[style=terminal_output]
Step  x       y         k1       k2       k3       k4       
------------------------------------------------------------
0     0.00    1.0000    0.0500   0.0525   0.0526   0.0551   
1     0.05    1.0525    0.0551   0.0578   0.0578   0.0605   
2     0.10    1.1103    0.0605   0.0633   0.0633   0.0662   
3     0.15    1.1737    0.0662   0.0691   0.0692   0.0721   
4     0.20    1.2428    

y(0.2): 1.24281
\end{lstlisting}
\end{tcolorbox}

\newpage

\chapter{Eigenvalue Problems \& Curve Fitting}

\section{Power Method for Dominant Eigenvalue}
\textbf{Category:} \textcolor{secondary}{Eigenvalues \& Curve Fitting} \quad | \quad \textbf{Code by:} \textcolor{secondary}{Anushka} \quad | \quad \textbf{Source:} \texttt{\detokenize{19_power_method_detailed.py}}

The Power Method is an iterative technique that computes the dominant eigenvalue (largest absolute magnitude) and its corresponding eigenvector for a square matrix $A$.

\begin{tcolorbox}[colback=theorybg,colframe=theoryframe,title=\textbf{Mathematical Theory \& Formulation},boxrule=1pt,arc=2mm]
\[
y_{k+1} = A x_k, \quad \lambda_{k+1} = \max(|y_{k+1}|), \quad x_{k+1} = \frac{y_{k+1}}{\lambda_{k+1}}
\]
\[
\text{Convergence Rate: } \left|\frac{\lambda_2}{\lambda_1}\right|^k \quad (\text{Requires } |\lambda_1| > |\lambda_2|)
\]
\end{tcolorbox}

\begin{tcolorbox}[colback=codebg,colframe=secondary,title=\textbf{Python Implementation (Step-by-Step Loops \& Pure Functions)},boxrule=1pt,arc=2mm]
\begin{lstlisting}[style=python_detailed]
# Code by Anushka

def power_method(A, tol=1e-5, max_iter=50):
    n = len(A)
    x = []
    for _ in range(n):
        x.append(1.0)
    lambda_old = 0.0

    print(f"{'Iter':<6}{'Eigenvalue':<14}{'Eigenvector':<20}")
    print("-" * 40)

    for step in range(1, max_iter + 1):
        # Matrix-vector multiplication y = A * x
        y = []
        for i in range(n):
            row_sum = 0.0
            for j in range(n):
                row_sum += A[i][j] * x[j]
            y.append(row_sum)

        # Find maximum absolute value in y
        lambda_new = y[0]
        for val in y:
            if abs(val) > abs(lambda_new):
                lambda_new = val

        # Normalize eigenvector x = y / lambda_new
        for i in range(n):
            x[i] = y[i] / lambda_new

        vec_str = "["
        for i in range(n):
            vec_str += f"{x[i]:.3f}"
            if i < n - 1:
                vec_str += ", "
        vec_str += "]"

        print(f"{step:<6}{lambda_new:<14.4f}{vec_str:<20}")

        if abs(lambda_new - lambda_old) < tol:
            return lambda_new, x

        lambda_old = lambda_new

    return lambda_new, x


matrix = [
    [4.0, 1.0],
    [2.0, 3.0]
]

val, vec = power_method(matrix)
rounded_vec = []
for v in vec:
    rounded_vec.append(round(v, 4))
print(f"\nEigenvalue: {val:.4f}, Eigenvector: {rounded_vec}")

\end{lstlisting}
\end{tcolorbox}

\begin{tcolorbox}[colback=termbg,colframe=primary,title=\textbf{Program Execution Output \& Convergence Table},boxrule=1pt,arc=2mm]
\begin{lstlisting}[style=terminal_output]
Iter  Eigenvalue    Eigenvector         
----------------------------------------
1     5.0000        [1.000, 1.000]      
2     5.0000        [1.000, 1.000]      

Eigenvalue: 5.0000, Eigenvector: [1.0, 1.0]
\end{lstlisting}
\end{tcolorbox}

\newpage

\section{Linear Curve Fitting (Least Squares)}
\textbf{Category:} \textcolor{secondary}{Eigenvalues \& Curve Fitting} \quad | \quad \textbf{Code by:} \textcolor{secondary}{Anushka} \quad | \quad \textbf{Source:} \texttt{\detokenize{20_linear_fit_detailed.py}}

Linear curve fitting finds the best-fit line $y = mx + c$ for a dataset of $n$ points by minimizing the sum of squared vertical residuals $\sum (y_i - (mx_i + c))^2$.

\begin{tcolorbox}[colback=theorybg,colframe=theoryframe,title=\textbf{Mathematical Theory \& Formulation},boxrule=1pt,arc=2mm]
\[
\text{Normal Equations: } n c + m \sum x = \sum y, \quad c \sum x + m \sum x^2 = \sum x y
\]
\[
m = \frac{n \sum xy - \sum x \sum y}{n \sum x^2 - (\sum x)^2}, \quad c = \frac{\sum y - m \sum x}{n}
\]
\end{tcolorbox}

\begin{tcolorbox}[colback=codebg,colframe=secondary,title=\textbf{Python Implementation (Step-by-Step Loops \& Pure Functions)},boxrule=1pt,arc=2mm]
\begin{lstlisting}[style=python_detailed]
# Code by Anushka

def linear_fit(x, y):
    n = len(x)

    # Compute summations using loops
    sx = 0.0
    sy = 0.0
    sxy = 0.0
    sx2 = 0.0

    for i in range(n):
        sx += x[i]
        sy += y[i]
        sxy += x[i] * y[i]
        sx2 += x[i] ** 2

    print(f"n = {n}, sum_x = {sx:.2f}, sum_y = {sy:.2f}, sum_xy = {sxy:.2f}, sum_x2 = {sx2:.2f}")

    m = (n * sxy - sx * sy) / (n * sx2 - sx**2)
    c = (sy - m * sx) / float(n)
    return m, c


x_data = [1.0, 2.0, 3.0, 4.0, 5.0]
y_data = [2.2, 2.8, 3.6, 4.5, 5.1]

m, c = linear_fit(x_data, y_data)
print(f"Fitted Line: y = {m:.4f}x + {c:.4f}")

\end{lstlisting}
\end{tcolorbox}

\begin{tcolorbox}[colback=termbg,colframe=primary,title=\textbf{Program Execution Output \& Convergence Table},boxrule=1pt,arc=2mm]
\begin{lstlisting}[style=terminal_output]
n = 5, sum_x = 15.00, sum_y = 18.20, sum_xy = 62.10, sum_x2 = 55.00
Fitted Line: y = 0.7500x + 1.3900
\end{lstlisting}
\end{tcolorbox}

\newpage

\chapter{Master Summary \& Quick-Reference Guide}

\begin{tcolorbox}[colback=codebg,colframe=secondary,title=\textbf{Unit 1: Roots of Non-Linear Equations},boxrule=1.2pt,arc=2mm]
\begin{itemize}[leftmargin=*,itemsep=4pt]
\item \textbf{1. Bisection Method:} $c = (a+b)/2.0$. Halves bracket interval. Order: Linear ($\mathcal{O}(2^{-n})$). Always converges if $f(a)f(b) < 0$.
\item \textbf{2. Regula-Falsi (False Position):} $c = \frac{a f(b) - b f(a)}{f(b) - f(a)}$. Secant line $x$-intercept with guaranteed bracketing. Order: Superlinear to Linear.
\item \textbf{3. Newton-Raphson Method:} $x_{k+1} = x_k - \frac{f(x_k)}{f'(x_k)}$. Tangent line root approximation. Order: Quadratic ($p=2$).
\item \textbf{4. Secant Method:} $x_{k+1} = x_k - f(x_k)\frac{x_k - x_{k-1}}{f(x_k) - f(x_{k-1})}$. Finite-difference slope approximation. Order: $p \approx 1.618$.
\item \textbf{5. Fixed-Point Iteration:} $x_{k+1} = g(x_k)$. Converges linearly if $|g'(x)| < 1$ in neighborhood of the fixed point.
\end{itemize}
\end{tcolorbox}

\begin{tcolorbox}[colback=codebg,colframe=secondary,title=\textbf{Unit 2: Systems of Linear Algebraic Equations},boxrule=1.2pt,arc=2mm]
\begin{itemize}[leftmargin=*,itemsep=4pt]
\item \textbf{6. Gauss Elimination:} Converts $[A \mid b]$ to upper-triangular matrix, followed by back substitution. Operations: $\mathcal{O}(\frac{2}{3}n^3)$.
\item \textbf{7. Gauss-Jordan Elimination:} Direct reduction to reduced row echelon form $[I \mid x]$ eliminating back substitution. Operations: $\mathcal{O}(n^3)$.
\item \textbf{8. Gauss-Jacobi Iteration:} $x_i^{(k+1)} = \frac{1}{a_{ii}}(b_i - \sum_{j \neq i} a_{ij} x_j^{(k)})$. Requires strict diagonal dominance.
\item \textbf{9. Gauss-Seidel Iteration:} $x_i^{(k+1)} = \frac{1}{a_{ii}}(b_i - \sum_{j < i} a_{ij} x_j^{(k+1)} - \sum_{j > i} a_{ij} x_j^{(k)})$. Converges $\approx 2\times$ faster than Jacobi.
\end{itemize}
\end{tcolorbox}

\begin{tcolorbox}[colback=codebg,colframe=secondary,title=\textbf{Unit 3: Interpolation \& Polynomial Approximation},boxrule=1.2pt,arc=2mm]
\begin{itemize}[leftmargin=*,itemsep=4pt]
\item \textbf{10. Lagrange Interpolation:} $P(x) = \sum y_i \prod_{j \neq i} \frac{x - x_j}{x_i - x_j}$. Flexible formulation for arbitrary, unequally spaced nodes.
\item \textbf{11. Newton's Divided Difference:} $P(x) = f[x_0] + \sum_{k=1}^n f[x_0,\dots,x_k] \prod_{j=0}^{k-1} (x - x_j)$. Ideal for incrementally adding data points.
\item \textbf{12. Newton's Forward Difference:} $P(x) = y_0 + u \Delta y_0 + \frac{u(u-1)}{2!} \Delta^2 y_0 + \dots$. Specifically for equispaced points ($u = (x-x_0)/h$).
\end{itemize}
\end{tcolorbox}

\begin{tcolorbox}[colback=codebg,colframe=secondary,title=\textbf{Unit 4: Numerical Quadrature \& Integration},boxrule=1.2pt,arc=2mm]
\begin{itemize}[leftmargin=*,itemsep=4pt]
\item \textbf{13. Trapezoidal Rule:} $I \approx \frac{h}{2}[f(a) + 2\sum_{i=1}^{n-1} f(x_i) + f(b)]$. Linear piecewise approximation. Error: $\mathcal{O}(h^2)$.
\item \textbf{14. Simpson's 1/3 Rule:} $I \approx \frac{h}{3}[f(a) + f(b) + 4\sum_{\text{odd}} + 2\sum_{\text{even}}]$. Parabolic piecewise fits. Requires even $n$. Error: $\mathcal{O}(h^4)$.
\item \textbf{15. Simpson's 3/8 Rule:} $I \approx \frac{3h}{8}[f(a) + f(b) + 3\sum_{i \nmid 3} + 2\sum_{i \mid 3}]$. Cubic piecewise fits. Requires $n$ multiple of 3. Error: $\mathcal{O}(h^4)$.
\end{itemize}
\end{tcolorbox}

\begin{tcolorbox}[colback=codebg,colframe=secondary,title=\textbf{Unit 5 \& 6: Differential Equations, Eigenvalues \& Curve Fitting},boxrule=1.2pt,arc=2mm]
\begin{itemize}[leftmargin=*,itemsep=4pt]
\item \textbf{16. Euler's Method:} $y_{n+1} = y_n + h f(x_n, y_n)$. Single-step explicit tangent method. Order: First-order $\mathcal{O}(h)$.
\item \textbf{17. Modified Euler (Heun's):} Predictor-corrector scheme. Predicts $\tilde{y}_{n+1}$, corrects using average slope. Order: Second-order $\mathcal{O}(h^2)$.
\item \textbf{18. Runge-Kutta 4 (RK4):} Four weighted intermediate slopes per step. Gold standard for initial value ODEs. Order: Fourth-order $\mathcal{O}(h^4)$.
\item \textbf{19. Power Method:} $y_{k+1} = A x_k, \ \lambda_{k+1} = \max(|y_{k+1}|), \ x_{k+1} = y_{k+1} / \lambda_{k+1}$. Iteratively converges to dominant eigenvalue.
\item \textbf{20. Linear Curve Fitting:} Minimizes $\sum (y_i - (m x_i + c))^2$ to determine optimal slope $m$ and intercept $c$ for best-fit straight line.
\end{itemize}
\end{tcolorbox}

\end{document}
